%-*-tex-*-
% Copyright Michael J. Ferguson, INRS-Telecommunications
% All rights reserved. 

%   ===== Citation generation ==========
% \cite is used with a journal/paper reference.
% The forms allow both random input order of cite forms or
% fixed order. At the moment it supports only IEEE form, ie 
% a list in the order of their occurrence in a paper. 
% A data base system using a mailbox would be much better. 


% Invocation of \cite{Ferg84} will define a \tag{Ferg84}{<number>},
% where the reference number is sequentially determined. At the site
% of the invocation <number> is dropped into the text (without [ or ]).
% If the reference has been previously invoked, the number is that of
% the first invocation.
% the \citenum is relative to an \i@nitcitenum. This allows for multiple
% reference lists in the same document. A new number is defined if the old
% one is undefined or less than or equal to the initcitenum

% In addition, when the citation tag file is written out, the last \citenum 
% is also written. This allows a citation tag file to be input for early 
% reference and then to be effectively discarded by resetting the
% \i@nitcitenum

\def\opencitationtagfile{\openlistfile{ctg}\resetcitationnumbers\immediate
             \write\ctg@file{\noexpand\initcitetag{\the\citenum}}}

\def\closecitationtagfile{\immediate
     \write\ctg@file{\noexpand\lastcitetag{\the\citenum}}\closelistfile{ctg}}

%--- citenum allocation ------
\def\initcitetag#1{\xdef\i@nitcitenum{#1}} 
                       \initcitetag{0} % initializes to 0
\def\lastcitetag#1{\xdef\l@astcitenum{#1}} 
                       \lastcitetag{0} % initializes to 0}

\newcount\citenum \citenum = 0  
\def\resetcitationnumbers{\ifnum\citenum>\l@astcitenum 
                            \lastcitetag{\the\citenum}\else
                            \global\citenum = \l@astcitenum \fi
                          \initcitetag{\l@astcitenum}}

\def\w@ritecite#1{\ifctglist \immediate\write\ctg@file
                         {\string\tag{#1}{\the\citenum}}\fi}
\def\m@akecitetag#1{\global\advance\citenum by 1 \relax
               \tag{#1}{\the\citenum}\w@ritecite{#1}}

\def\cite#1{\ifctglist \else \opencitationtagfile \fi
            \edef\c@itearrow{<--}\relax
            \ifundefined{:@#1}\m@akecitetag{#1}\edef\c@itearrow
                    {==>}\fi
            \ifnum \i@nitcitenum < \csname :@#1\endcsname \else
                   \m@akecitetag{#1}\edef\c@itearrow
                    {==>}\fi
             \count255= \csname :@#1\endcsname 
             \advance\count255 by -\i@nitcitenum \relax
             \proofmargin{#1 \c@itearrow\ \the\count255}\the\count255}

% This will produce the correct value of a citation without 
% causing it to be created if it does not exist

\def\citeref#1{\ifundefined{:@#1}\count255=0\else
                     \count255=\csname :@#1\endcsname\fi
          \advance\count255by-\i@nitcitenum\relax\the\count255}


% Citation/Reference list generation. An example setup using the list
% macros to generate a reference list is shown below.
%    \opencitationtagfile
%      <main text>
%    \beginlist
%      \makecitationlist %or \makelongcitationlist
%    \endlist
% The reference list is generated from two files:
%    <jobname>.ctg  contains citation tags of the form \tag{Ferg84}{4}
%                   where the first argument is the tag, and the second
%                   is the reference number associated with that tag.
%    <jobname>.cfm  contains the text for the references. These can appear
%                   in any order and contain text of the form
%                   \citeform{<tag>}{\listitem{\citetagvalue} <reference text>}
%                   This example shows a preamble to the text which invokes
%                   a new list item with the appropriate reference number.
%                   The reference text is the author, journal etc. for the
%                   reference.
% Each \citeform entry will generate a new macro definition of the form
% \def\@Ferg84form{<reference text>}.
% Each \tag entry will define \citetagvalue as the citation/reference number
% and invoke \@Ferg84form 

%\makecitationlist does not require the .cfm file to have citations in the 
%same order as .ctg file. \makelongcitationlist enters the citations in the
%in the document in the same order as in the file. It numbers with a counter
%and indicates if there is a difference between citetag and counter.

\newif\ifctglist  % so dont have to check if undefined
\newtoks\citetagfilename  \citetagfilename = {\jobname.ctg} %compatible Ugh!
\newtoks\inputciteformfiles 
       \inputciteformfiles = {\inputwithcheck{\jobname.cfm}}
\newtoks\inputcitetagfiles 
   \inputcitetagfiles = {\inputwithcheck{\the\citetagfilename}}



%\def\citetagfileclosedmessage{ << CITATION TAG FILE CLOSED >> }
%\def\nocitetagformmessage{ NO CITATION FORM FOR THIS TAG}
%            see english/french titles

\def\c@losectg{
             \ifctglist 
               \immediate\write\ctg@file{\noexpand\lastcitetag{\the\citenum}}
               \immediate\closeout\ctg@file
               \global\ctglistfalse
               \writeterm{\citetagfileclosedmessage}{}{}\relax
             \fi}

% outside so that the cfm list may be included in the document.
\def\citeform#1#2{\expandafter
                      \def\csname @#1form\endcsname{#2}}


%cfm entries may be in any order
\long\def\makecitationlist{\par
           \begingroup 
             \c@losectg
             \def\citeform##1##2{\expandafter
                      \def\csname @##1form\endcsname{##2}}
             \def\tag##1##2{\count255 = ##2\relax
                            \advance\count255 by -\i@nitcitenum
                               \edef\citetagvalue{\the\count255}\relax  % locally redefine \tag
                            \ifundefined{@##1form}\relax
                              \writeterm{<<<\the\count255. \nocitetagformmessage: ##1}{}{}\relax
                              \vskip 2ex
                              \leftline{##2\  \nocitetagformmessage: ##1}
                              \vskip 2ex
                            \else
                              \csname @##1form\endcsname
                            \fi}
             \the\inputciteformfiles
             \the\inputcitetagfiles \par
           \endgroup}

%.cfm entries must be in the same order as the tags in the  .ctg file
\newcount\c@fmcount %counter for citations
\long\def\makelongcitationlist{\par
            \begingroup
               \c@losectg
               \c@fmcount=0 \relax 
               \def\citetagvalue{\the\c@fmcount}
               \def\citeform##1##2{\advance\c@fmcount by 1 
                                   \ifundefined{:@##1} 
                                       \expandafter \xdef \csname :@##1\endcsname{-1} \fi
                                   \count 255 = \csname :@##1\endcsname 
                                   \advance\count 255 by -\i@nitcitenum 
                                   \ifnum \the\count255 = \the\c@fmcount
                                      \relax \else
         \writeterm{<<< Cite tag error -- Tag ##1 -- Citation \the\c@fmcount >>>}{}{}
                                     \fi
                                   ##2}
              \the\inputcitetagfiles 
              \the\inputciteformfiles\par
             \endgroup
            }

